\section{Fertigung, Bestückung und Inbetriebnahme}
Nach Abschluss des Layouts wurden Fertigungsdaten erzeugt, die Platine bestellt und bestückt. In der ersten Inbetriebnahme stand die grundlegende elektrische Funktionsprüfung im Vordergrund. Ziel dieser Phase war es, die Versorgung, die Erreichbarkeit der zentralen Komponenten und offensichtliche Hardwarefehler früh zu prüfen.

Die folgende Tabelle fasst den aktuellen Stand der Inbetriebnahme zusammen. Dabei handelt es sich um Funktionsprüfungen und grobe Spannungsprüfungen, noch nicht um eine vollständige messtechnische Charakterisierung.

\begin{table}[H]
	\centering
	\small
	\renewcommand{\arraystretch}{1.35}
	\begin{tabularx}{\linewidth}{>{\RaggedRight\arraybackslash}m{0.22\linewidth}|>{\RaggedRight\arraybackslash}m{0.29\linewidth}|>{\RaggedRight\arraybackslash}m{0.25\linewidth}|>{\Centering\arraybackslash}m{0.11\linewidth}}
		\hline
		\textbf{Teilbereich} & \textbf{Prüfung} & \textbf{Ergebnis} & \textbf{Status} \\\hline
		Spannungsversorgung & Kontrolle der erzeugten Versorgungsschienen & 6-V-Bereiche liefern ca. 6 V, Raspberry-Pi-Schiene ca. 5 V, Logikversorgung ca. 3,3 V & funktioniert \\\hline
		STM32 & Programmieren und Erreichen des Mikrocontrollers über SWD & STM32 kann geflasht werden & funktioniert \\\hline
		Motoransteuerung & STM32 setzt PWM-Signale an den Motortreibern & Motortreiber reagieren auf die vom STM32 ausgegebenen Steuersignale & funktioniert \\\hline
		UART-Kommunikation & Steuerbefehle vom Raspberry Pi beziehungsweise Testsystem an den STM32 & UART-Kommandos können verarbeitet und in PWM-Werte umgesetzt werden & funktioniert \\\hline
		IMU LSM6DSO & Zugriff vom Raspberry Pi über die verwendete Softwarebibliothek & Sensor ist erreichbar und kann angesprochen werden & funktioniert \\\hline
		IMU BNO085 & Zugriff vom Raspberry Pi über die verwendete Softwarebibliothek & Sensor ist derzeit nicht erreichbar & offen \\\hline
		Servoansteuerung & Ansteuerung des PCA9685 und Prüfung des PWM-Ausgangs & Gewünschtes PWM-Signal wird erzeugt & funktioniert \\\hline
	\end{tabularx}
	\caption{Aktueller Stand der Inbetriebnahme}
\end{table}


\section{Fehler, Probleme und bereits umgesetzte Verbesserungen}
\subsection{Fehler im Schaltplan aus dem Power-Designer}
Ein konkreter Fehler trat bei der Übernahme eines Schaltungsteils aus dem Texas-Instruments-Power-Designer auf. Beim Abzeichnen der erzeugten Schematic wurde ein Fehler übernommen, der erst im späteren Verlauf erkannt wurde.

Die Korrektur erfolgte direkt an der aufgebauten Platine durch das Auftrennen einer Leiterbahn. Damit konnte der Fehler behoben werden. Für eine nächste Platinenrevision muss diese Änderung jedoch sauber in Schaltplan und Layout zurückgeführt werden.

\subsection{Empfindlichkeit des 5-V-Wandlers beim Abziehen des Akkus}
Ein weiteres Problem betrifft den 5-V-Wandler für den Raspberry Pi. Wenn der Akku während des Betriebs einfach abgezogen wird, kann der Wandler beschädigt werden. Deshalb darf der Akku derzeit nur getrennt werden, nachdem der Raspberry Pi sauber heruntergefahren und abgeschaltet wurde.

Für eine zukünftige Revision sollte daher geprüft werden, ob zusätzliche Schutzmaßnahmen, ein geänderter Abschaltpfad oder eine robustere Wandlerlösung notwendig sind.


